\addtocontents{toc}{\protect\Needspace{3\baselineskip}}
\originalchapter{二形式规范对称与开弦端点电荷}\label{ch:basic-twoform}
对应 8.251 第 25 讲. 此讲的关键是开弦端点如何与闭弦二形式的规范对称相容. 直接写一个二形式体耦合会留下边界项, 膜上的规范场恰能补偿它. 原课笔记的系数、指标有若干混写, 本稿采用以下一致约定, 并保留其物理构造路线.

\originalsection{二形式及带取向世界面耦合}
令 $B=\frac12B_{\mu\nu}\dd x^\mu\wedge\dd x^\nu$, $H=\dd B$, 因而
\[
 H_{\mu\nu\rho}=\partial_\mu B_{\nu\rho}
 +\partial_\nu B_{\rho\mu}+\partial_\rho B_{\mu\nu},
 \qquad B\mapsto B+\dd\Lambda.
\]
为把 $2\pi\alpha'$ 系数清楚分开, 先用已将系数吸收进 $B$ 的约定, 取
\[
 S_B=\int_\Sigma X^*B
 =\frac12\int\dd\tau\dd\sigma\,
 \epsilon^{ab}B_{\mu\nu}(X)\partial_aX^\mu\partial_bX^\nu,
 \quad\epsilon^{\tau\sigma}=1.
\]
取向反转使此耦合变号, 正是二形式电荷对弦取向敏感的原因. $H$ 在 $B\mapsto B+\dd\Lambda$ 下不变, 因 $\dd^2=0$.

\originalsection{弦电流及其端点散度}
以 $S_B=\frac12\int\dd^Dx B_{\mu\nu}J^{\mu\nu}$ 定义
\[
 J^{\mu\nu}(x)=\int_\Sigma\dd\tau\dd\sigma\,
 (\dot X^\mu X'^\nu-\dot X^\nu X'^\mu)
 \delta^{(D)}(x-X).
\]
对光滑紧支撑测试一形式 $\Lambda$, 积分分部与世界面 Stokes 定理给
\[
 -\int\dd^Dx\,\Lambda_\nu\partial_\mu J^{\mu\nu}
 =\int_{\partial\Sigma}\Lambda_\nu(X)\dd X^\nu.
\]
所以闭弦无真实边界时 $\partial_\mu J^{\mu\nu}=0$, 开弦则在两个端点留下带相反符号的分布源. 这个测试函数推导避免对 $\delta(x-X)$ 直接猜符号. 实际初末时间切片也有边界, 通常令变换在那里为零, 单独关注空间端点.

静态规范 $X^0=t$ 下,
\[
 J^{0i}(t,\boldsymbol x)=\int\dd\sigma\,
 X'^i\delta^{(d)}(\boldsymbol x-\boldsymbol X(t,\sigma)).
\]
它沿弦的正取向. 对静止闭弦, 空间散度为零, 与稳恒电流的形状相似. 对开弦, 散度等于端点的正负 $\delta$ 分布, 说明弦电荷的流线不能无补偿地突然终止.

\originalsection{三维空间中的稳恒电流类比}
为固定场方程系数, 取
\[
 S_H=-\frac1{12\kappa_B^2}\int\dd^Dx\,H_{\mu\nu\rho}H^{\mu\nu\rho}
 +\frac12\int\dd^Dx\,B_{\mu\nu}J^{\mu\nu}.
\]
变分并用 $H$ 的反对称性得 $\partial_\rho H^{\rho\mu\nu}=-\kappa_B^2J^{\mu\nu}$. 再求 $\partial_\mu$ 自动要求体源散度为零, 与闭弦电流相容. 不同的耦合正负号会同时改变源的定义, 物理结论不变.

在 $3+1$ 维的静止情形可将 $H^{i0j}$ 写成空间 $\epsilon^{ijk}$ 与辅助矢量 $\mathcal B_k$ 的对偶, 场方程成为 $\nabla\times\boldsymbol{\mathcal B}=\boldsymbol J^0$, 选择对偶符号与上式相容. 这是求闭弦二形式场与稳恒电流磁场的数学类比. $B_{\mu\nu}$ 对应势、$H$ 对应场强, 不能因同用字母 $B$ 就把二形式势直接等同于普通三维磁场.

\originalsection{规范变换留下了什么}
有 $\delta B=\dd\Lambda$ 时
\[
 \delta S_B=\int_\Sigma X^*(\dd\Lambda)
 =\int_{\partial\Sigma}X^*\Lambda.
\]
闭弦空间端点粘合, 对应贡献相消. 开弦的两个边界世界线带相反的诱导取向, 一般不为零. 若端点落在 D 膜, 固定法向坐标使 $\dd X^a=0$, 但膜内 $\Lambda_m\dd X^m$ 仍然存在. 所以只将法向固定不能解决二形式规范变换的边界问题.

\originalsection{端点规范场如何完成补偿}
在膜上加入一形式 $A=A_m\dd x^m$ 的端点耦合
\[
 S_{\partial}=\int_{\partial\Sigma}A,
 \qquad\delta A=-\Lambda|_W.
\]
于是 $\delta(S_B+S_{\partial})=0$. 两个端点的相反取向给相反电荷, 与粒子、反粒子式端点标签一致. 原课所说“弦无法在膜上终止”的中间障碍, 因而被一个具体的边界自由度解决, 不必丢掉二形式规范对称.

但 $F=\dd A$ 现在也变化为 $\delta F=-\dd\Lambda|_W$. 在吸收系数的约定中, 不变量是 $\mathcal F=B|_W+F$. 恢复通常弦论中几何二形式 $B^{\rm phys}$ 的归一化,
\[
 \mathcal F=B^{\rm phys}|_W+2\pi\alpha'F,
 \qquad\delta A=-\frac{\Lambda^{\rm phys}|_W}{2\pi\alpha'}.
\]
同时仍有普通 $A\mapsto A+\dd\epsilon$ 的膜上 $U(1)$ 变换. 二者分别补偿体二形式与膜上规范冗余, 不应把它们当同一个标量参数变换.

\originalsection{膜上场强承接弦电荷}
膜作用量必须依 $\mathcal F$ 而非单独 $F$, 最低阶结构为 $-\frac1{4g^2}\int_W\mathcal F_{mn}\mathcal F^{mn}$. 展开含 $F^2,B^2,2F\cdot B$. 因此膜上电场也为二形式提供相应源, 可在端点外继续承接弦电荷流. 它不表示基本弦严格由普通电场线组成, 而是体与边界源通过规范对称相联系.

在完整 DBI 作用量中同样出现 $\sqrt{-\det(G+\mathcal F)}$, 第二编将展开其张力与规范耦合. 此处已完整解释所需的边界取消机制, 无需在入门章重复全息编的非阿贝尔低能理论.

\originalsection{原课末尾模型图的适用范围}
第 25 讲用不同颜色膜、弱相互作用膜和环面上的相交膜示意标准模型构造: 端点标签可承载不同规范表示, 膜间弦可能形成手征物质. 这些图表达模型构建路线, 不是一套已在本讲检查了异常、模稳定、对称破缺和所有实验约束的实际标准模型. 本稿不转载其中的第三方示意图, 也不把图中的每种端点命名当作已得到真实夸克与轻子全部质量的计算结果. 精确的相交膜紧致化构造保留为待补方向.
